\documentclass[11pt,a4paper]{amsart}

\usepackage[utf8]{inputenc}
\usepackage[T1]{fontenc}
\usepackage{lmodern}
\usepackage[polish]{babel}
\usepackage{amsmath, amssymb, amsthm}
\usepackage{hyperref}
\usepackage{enumerate}
\usepackage{geometry}
\usepackage{float}
\usepackage[section]{placeins}
\usepackage{longtable}
\usepackage{booktabs}
\usepackage{csquotes}

\theoremstyle{plain}
\newtheorem{theorem}{Twierdzenie}[section]
\newtheorem{lemma}[theorem]{Lemat}
\newtheorem{corollary}[theorem]{Wniosek}
\newtheorem{conjecture}[theorem]{Hipoteza}

\theoremstyle{definition}
\newtheorem{definition}[theorem]{Definicja}

\theoremstyle{remark}
\newtheorem{remark}[theorem]{Uwaga}

\usepackage[
    backend=biber,
    style=numeric,
    sorting=none
]{biblatex}
\addbibresource{bibliography.bib}

\geometry{margin=2.8cm}

\title{Liniowy algorytm formy kwadratowej Huanga\\
na półgrupach numerycznych i asymptotyka gęstości}

\author{Artur Flamandzki}
\thanks{Adres e-mail: \texttt{flamandzki.artur@gmail.com}.}

\date{\today}

\begin{document}

\begin{abstract}
Niech $G = \mathbb{N} \setminus \langle a, b \rangle$ będzie zbiorem luk
półgrupy numerycznej dla wzajemnie pierwszych $a < b$, oraz $N = |G|$.
Huang~\cite{Huang2026} zdefiniował formę kwadratową
$Q(\mathbf{n}) = \sum_{i,j} K(j-i)\,n_i n_j$ na $\mathbb{R}^G$, gdzie
$K(d) = \mathbf{1}_{d \geq 0} - \mathbf{1}_{d \geq a} - \mathbf{1}_{d \geq b}
+ \mathbf{1}_{d \geq a+b}$, i wykazał, że odtwarza ona statystykę
$\mathrm{dinv}$ na wymiernych ścieżkach Dycka.
Naiwne obliczenie $Q$ wymaga $O(N^2)$ operacji.

Dowodzimy, że przedziałowa struktura $K$ pozwala zredukować $Q$ do
liniowej kombinacji sum z przesuwnym oknem, dając algorytm czasu $O(N)$
(Twierdzenie~\ref{thm:linear}). Uogólniamy na półgrupy
$\langle p_1,\ldots,p_n \rangle$ o $n$ generatorach, gdzie uogólnione
jądro $K^{(n)}$ zdefiniowane zasadą włączeń--wyłączeń
--- o $2^n$ składnikach --- posiada co najwyżej $w(n) \leq \sigma_n$
okien aktywnych wskutek kasowania parzystościowego;
są one obliczalne inkrementalnie w czasie $O(n \cdot \sigma_n)$
(Twierdzenie~\ref{thm:general}).
Empirycznie $w(n) \sim \sigma_n$ przy $w(n) \ll 2^n$:
wykładnicza kombinatoryka redukuje się do liniowej struktury okien.
Dla ciągu parami różnych liczb całkowitych $p_i \geq 2$ spełniającego
$\log p_n = o(n)$ dowodzimy
\[
  \frac{w(n)}{2^n - 1} \;=\; o\!\left(\,\prod_{i=1}^{n}
    \Bigl(1 - \tfrac{1}{p_i}\Bigr)\right), \qquad n \to \infty
\]
(Twierdzenie~\ref{thm:density}).
\end{abstract}

\maketitle

% ==================================================
\section{Wprowadzenie}
\label{sec:intro}
% ==================================================

\subsection*{1.1. Kontekst i motywacja}

Niech $a, b \in \mathbb{N}$, $\gcd(a,b) = 1$, $a < b$.
Półgrupa numeryczna generowana przez $a$ i $b$ to
$\langle a, b \rangle = \{ xa + yb : x, y \in \mathbb{Z}_{\geq 0} \}$.
Na mocy twierdzenia Sylvestera--Frobeniusa zbiór luk
$G = \mathbb{N} \setminus \langle a, b \rangle$ jest skończony,
$|G| = N = \tfrac{(a-1)(b-1)}{2}$.
Piszemy $G = \{g_0 < g_1 < \cdots < g_{N-1}\}$.

Huang~\cite{Huang2026} zdefiniował na $\mathbb{R}^G$ formę kwadratową
\[
  Q(\mathbf{n}) = \sum_{i,j \in G} K(j - i)\, n_i n_j,
  \qquad
  K(d) = \mathbf{1}_{d \geq 0} - \mathbf{1}_{d \geq a}
         - \mathbf{1}_{d \geq b} + \mathbf{1}_{d \geq a+b},
\]
i wykazał, że gdy $\mathbf{n} = \mathbf{1}_D$ jest wektorem wskaźnikowym
poddiagramu Younga $D \subseteq G$, to $Q(\mathbf{1}_D) = \mathrm{dinv}(D)$,
gdzie $\mathrm{dinv}$ jest statystyką Gorsky'ego--Mazina na wymiernych
ścieżkach Dycka. Niezależna formalna weryfikacja tego twierdzenia
w systemie Lean~4/Mathlib: \cite{AxiomProver2026}.

Praca~\cite{Huang2026} oraz jej formalizacja~\cite{AxiomProver2026}
dotyczą tożsamości $Q(\mathbf{1}_D) = \mathrm{dinv}(D)$.
Niniejsza praca dotyczy efektywności obliczeniowej formy $Q$
i jej uogólnienia na $n$ generatorów.

Bezpośrednia ewaluacja $Q$ ze wzoru definicyjnego wymaga
$N^2$ operacji. Kluczową obserwacją jest, że $K$ przyjmuje
co najwyżej trzy wartości niezerowe na pięciu rozłącznych
przedziałach (Lemat~\ref{lem:K-intervals}), co redukuje
obliczenie $Q$ do dwóch sum z przesuwnym oknem, dając
algorytm $O(N)$.

Generalizacja na $n$ generatorów definiuje jądro $K^{(n)}$
zasadą włączeń--wyłączeń po wszystkich $2^n$ podzbiorach
$\{p_1, \ldots, p_n\}$. Struktura przedziałowa $K^{(n)}$
zachowuje się, lecz liczba okien aktywnych $w(n)$
--- interwałów, na których $K^{(n)} \neq 0$ ---
jest znacznie mniejsza niż $2^n$ wskutek kasowania
parzystościowego. Udowadniamy $w(n) \leq \sigma_n$
i podajemy algorytm inkrementalny obliczający $w(n)$
i same okna aktywne.

Stosunek $w(n)/(2^n - 1)$ ma naturalne powiązanie z iloczynem
$P(n) = \prod_{i=1}^n (1 - 1/p_i)$: przy warunku
$\log p_n = o(n)$ wykazujemy $w(n)/(2^n-1) = o(P(n))$,
tj.\ gęstość okien aktywnych jest asymptotycznie zaniedbywalnie
mała względem $P(n)$.

Uogólnienie na $n$ generatorów nie dotyczy problemu Frobeniusa,
który pyta o wzór jawny na największy element
$\mathbb{Z}_{\geq 0} \setminus \langle p_1,\ldots,p_n \rangle$
i pozostaje otwarty dla $n \geq 3$~\cite{Ramirez2005}.
Jądro $K^{(n)}$ jest zdefiniowane przez zasadę włączeń--wyłączeń
nad zbiorem generatorów; algorytm przyjmuje $G$ jako wejście.

\subsection*{1.2. Definicje i główne wyniki}

\textbf{Dwa generatory.}
Dla ustalonego $k$ i $\mathbf{n} = (n_j)_{j \in G} \in \mathbb{R}^G$
definiujemy sumy okienkowe
\begin{align*}
  W^+_k &:= \sum_{\substack{j \in G \\ g_k \leq g_j < g_k + a}} n_j, \qquad
  W^-_k := \sum_{\substack{j \in G \\ g_k + b \leq g_j < g_k + a + b}} n_j.
\end{align*}

\begin{theorem}[Algorytm liniowy, dwa generatory]
\label{thm:linear}
Dla każdego $\mathbf{n} \in \mathbb{R}^G$,
\[
  Q(\mathbf{n}) = \sum_{k=0}^{N-1} n_k \bigl(W^+_k - W^-_k\bigr).
\]
Obie rodziny $(W^+_k)$ i $(W^-_k)$ są obliczalne łącznie
w czasie $O(N)$ i pamięci $O(N)$ za pomocą czterech
wskaźników przesuwnego okna na $G$.
\end{theorem}

Twierdzenie~\ref{thm:linear} redukuje koszt z $O(N^2)$ do $O(N)$.
Ponieważ $N = \tfrac{(a-1)(b-1)}{2}$, algorytm działa w czasie
$O(ab)$ wobec $O(a^2 b^2)$ dla naiwnego.

\textbf{Wiele generatorów.}
Niech $\mathbf{p} = (p_1, \ldots, p_n)$, $p_i \geq 2$ parami różne,
$\gcd(p_1,\ldots,p_n) = 1$, $\sigma_n = \sum_{i=1}^n p_i$.
Uogólnione jądro
\[
  K^{(n)}(d) := \sum_{S \subseteq \{p_1,\ldots,p_n\}} (-1)^{|S|}
               \mathbf{1}_{d \geq \sigma(S)}, \qquad
  \sigma(S) = \sum_{p \in S} p, \quad \sigma(\emptyset) = 0,
\]
jest niezerowe co najwyżej na $w(n)$ przedziałach między kolejnymi
punktami podziału; piszemy $w(n)$ na liczbę tych okien aktywnych.

\begin{theorem}[Algorytm z przesuwnym oknem, $n$ generatorów]
\label{thm:general}
Niech $G = \mathbb{N} \setminus \langle p_1,\ldots,p_n \rangle$.
Forma $Q(\mathbf{n}) = \sum_{i,j \in G} K^{(n)}(j-i)\,n_i n_j$
jest obliczalna w czasie $O(T_{\mathrm{prep}} + w(n) \cdot N)$, gdzie
$T_{\mathrm{prep}} = O(n \cdot \sigma_n)$.
Zachodzi $w(n) \leq \sigma_n$.
Zbiór okien aktywnych i wartości $K^{(n)}$ na nich są obliczalne
inkrementalnie: przejście od $n-1$ do $n$ generatorów kosztuje
$O(\sigma_{n-1})$.
\end{theorem}

\textbf{Asymptotyka gęstości.}
Dla ciągu parami różnych liczb całkowitych $p_i \geq 2$
definiujemy iloczyn gęstości i deficyt gęstości:
\[
  P(n) := \prod_{i=1}^{n} \!\left(1 - \frac{1}{p_i}\right), \qquad
  \delta(n) := P(n) - \frac{w(n)}{2^n - 1}.
\]

\begin{remark}
\label{rem:terminology}
Termin \emph{luka} oznacza wyłącznie element
$G = \mathbb{N} \setminus \langle p_1,\ldots,p_n \rangle$.
Symbol $\delta(n)$ oznacza wyłącznie wielkość $P(n) - w(n)/(2^n-1)$.
Oba pojęcia są niezależne.
\end{remark}

\begin{remark}
\label{rem:mertens}
Iloczyn $P(n) = \prod_{i=1}^n (1 - 1/p_i)$ jest zdefiniowany
dla dowolnego ciągu parami różnych liczb całkowitych $p_i \geq 2$
i nie wymaga, by $p_i$ były liczbami pierwszymi.
Gdy $p_i$ są kolejnymi liczbami pierwszymi, $P(n)$ pokrywa się
z klasycznym iloczynem Mertensa; w tym przypadku $P(n) \to 0$.
\end{remark}

\begin{theorem}[Asymptotyka gęstości]
\label{thm:density}
Niech $(p_i)_{i \geq 1}$ będzie ciągiem parami różnych
liczb całkowitych $\geq 2$ spełniającym $\log p_n = o(n)$.
Wówczas
\[
  \frac{w(n)}{2^n - 1} = o(P(n)), \qquad n \to \infty,
\]
równoważnie: $\delta(n)/P(n) \to 1$.
\end{theorem}

Warunek $\log p_n = o(n)$ jest spełniony przez każdą rodzinę
generatorów pierwszych ($p_n \sim n \log n$), liczby pierwsze
bliźniacze, postępy arytmetyczne liczb pierwszych oraz ciągi
losowych liczb nieparzystych o gęstości pierwszej.

% ==================================================
\section{Preliminaria}
\label{sec:prelim}
% ==================================================

\subsection*{Jądro dwugeneratorowe}

\begin{lemma}[Struktura przedziałowa $K$]
\label{lem:K-intervals}
Dla wszystkich $d \in \mathbb{Z}$:
\[
  K(d) =
  \begin{cases}
     0  & d < 0,            \\
    +1  & 0 \leq d < a,     \\
     0  & a \leq d < b,     \\
    -1  & b \leq d < a+b,   \\
     0  & d \geq a+b.
  \end{cases}
\]
\end{lemma}

\begin{proof}
Bezpośrednie zestawienie czterech wskaźników $\mathbf{1}_{d \geq 0}$,
$-\mathbf{1}_{d \geq a}$, $-\mathbf{1}_{d \geq b}$,
$+\mathbf{1}_{d \geq a+b}$ przy $a < b$.
Pięć przypadków jest wzajemnie wykluczające się i wyczerpujące.
\end{proof}

\begin{corollary}
\label{cor:support}
$K(d) \neq 0$ implikuje $0 \leq d < a+b$.
Suma $\sum_{d \in \mathbb{Z}} K(d) = a - (b - a) = 2a - b$.
\end{corollary}

\subsection*{Jądro wielogeneratorowe}

\begin{definition}[Funkcja parytetu]
\label{def:parity}
Dla $\mathbf{p} = (p_1,\ldots,p_n)$ i $s \in \mathbb{Z}_{\geq 0}$:
\[
  \pi_n(s) := \sum_{\substack{S \subseteq \{p_1,\ldots,p_n\} \\
                              \sigma(S) = s}} (-1)^{|S|}.
\]
Zachodzi $K^{(n)}(d) = \sum_{s \leq d} \pi_n(s)$,
a zatem $K^{(n)}$ jest stałe na przedziałach między kolejnymi
$s$ z $\pi_n(s) \neq 0$.
\end{definition}

\begin{lemma}[Znikanie poza $[0, \sigma_n)$]
\label{lem:K-vanish}
$K^{(n)}(d) = 0$ dla $d < 0$ oraz dla $d \geq \sigma_n$.
\end{lemma}

\begin{proof}
Dla $d < 0$: każdy wskaźnik $\mathbf{1}_{d \geq \sigma(S)} = 0$.
Dla $d \geq \sigma_n$: wszystkie wskaźniki równe $1$, i
$\sum_{S \subseteq \{p_1,\ldots,p_n\}} (-1)^{|S|} = (1-1)^n = 0$
na mocy wzoru dwumianowego.
\end{proof}

\begin{lemma}[Rekurencja parytetowa]
\label{lem:parity-recurrence}
\[
  \pi_{n+1}(s) = \pi_n(s) - \pi_n(s - p_{n+1}).
\]
\end{lemma}

\begin{proof}
Podzbiory $\{p_1,\ldots,p_{n+1}\}$ o sumie $s$ dzielą się na:
podzbiory $S \subseteq \{p_1,\ldots,p_n\}$ z $\sigma(S) = s$,
wnoszące $\pi_n(s)$, oraz podzbiory $S \cup \{p_{n+1}\}$
z $\sigma(S) = s - p_{n+1}$, wnoszące $-\pi_n(s - p_{n+1})$.
\end{proof}

% ==================================================
\section{Dowód Twierdzenia~\ref{thm:linear}}
\label{sec:proof-linear}
% ==================================================

\subsection*{3.1. Redukcja formy kwadratowej}

\begin{lemma}
\label{lem:inner-sum}
Dla ustalonego $i \in G$:
\[
  \sum_{j \in G} K(j - i)\, n_j = W^+_i - W^-_i.
\]
\end{lemma}

\begin{proof}
Z Lematu~\ref{lem:K-intervals}: $K(j - i) = +1$ gdy
$g_i \leq g_j < g_i + a$, $K(j-i) = -1$ gdy
$g_i + b \leq g_j < g_i + a + b$, i $K(j-i) = 0$ w pozostałych
przypadkach. Definicje $W^+_i$ i $W^-_i$ odpowiadają dokładnie
tym dwóm przedziałom.
\end{proof}

\begin{proof}[Dowód Twierdzenia~\ref{thm:linear}]
Przestawiając sumę:
\[
  Q(\mathbf{n}) = \sum_{i \in G} n_i \sum_{j \in G} K(j-i)\,n_j
  = \sum_{i \in G} n_i (W^+_i - W^-_i),
\]
gdzie pierwsza równość wynika z symetrii (zamiana indeksów),
a druga z Lematu~\ref{lem:inner-sum}.

Liniowość: utrzymujemy cztery wskaźniki
$\mathrm{lo}^+, \mathrm{hi}^+, \mathrm{lo}^-, \mathrm{hi}^-$
do posortowanego $G$, zainicjalizowane na $0$.
Dla $k = 0, 1, \ldots, N-1$:
przesuwamy $\mathrm{hi}^+$ dopóki $g_{\mathrm{hi}^+} < g_k + a$;
przesuwamy $\mathrm{lo}^+$ dopóki $g_{\mathrm{lo}^+} < g_k$;
przesuwamy $\mathrm{hi}^-$ dopóki $g_{\mathrm{hi}^-} < g_k + a + b$;
przesuwamy $\mathrm{lo}^-$ dopóki $g_{\mathrm{lo}^-} < g_k + b$.
Każdy wskaźnik jest niemalejący i przekracza $G$ co najwyżej raz:
łącznie $\leq 4N$ aktualizacji.
\end{proof}

\begin{corollary}
\label{cor:speedup}
Dla $(a, b) = (97, 101)$, $N = 4801$: algorytm $O(N)$ wykonuje
$19\,204$ operacji wobec $23\,049\,601$ dla $O(N^2)$;
stosunek $\approx 1200$.
Ogólnie $N^2 / (4N) = N/4 = (a-1)(b-1)/8$.
\end{corollary}

% ==================================================
\section{Dowód Twierdzenia~\ref{thm:general}}
\label{sec:proof-general}
% ==================================================

\subsection*{4.1. Okna aktywne i ograniczenie $w(n) \leq \sigma_n$}

\begin{definition}[Okno aktywne]
\label{def:active}
Przedział $[s, s')$ między kolejnymi punktami podziału $K^{(n)}$
jest \emph{aktywny}, jeśli $K^{(n)} \neq 0$ na $(s, s')$.
$w(n)$ oznacza liczbę okien aktywnych.
\end{definition}

\begin{lemma}
\label{lem:w-bound}
$w(n) \leq \sigma_n$.
\end{lemma}

\begin{proof}
Z Lematu~\ref{lem:K-vanish} nośnik $K^{(n)}$ zawiera się
w $[0, \sigma_n) \cap \mathbb{Z}$, który zawiera $\sigma_n$
liczb całkowitych. Liczba okien aktywnych nie przekracza
liczby punktów podziału, co jest co najwyżej $\sigma_n$.
\end{proof}

\begin{lemma}[Kasowanie parzystościowe]
\label{lem:cancel}
Jeśli $\pi_n(s) = 0$, to punkt $s$ nie jest punktem podziału
$K^{(n)}$, co zmniejsza $w(n)$ poniżej $\sigma_n$.
Efekt ten zachodzi, gdy istnieją $S_1 \neq S_2$ z
$\sigma(S_1) = \sigma(S_2) = s$ i $|S_1| \not\equiv |S_2| \pmod{2}$.
\end{lemma}

\begin{proof}
Z definicji $\pi_n$: wkłady z przeciwną parzystością anulują się.
\end{proof}

\begin{remark}
\label{rem:cancellation-empirical}
Współczynnik nadmiarowości
$1 - w(n)/\sigma_n$ wzrasta od $0{,}0\%$ przy $n=2$
do $100{,}0\%$ przy $n \geq 21$ dla ciągu pierwszych liczb
pierwszych (Tabela~\ref{tab:mertens};
wartości $100{,}0\%$ wynikają z zaokrąglenia
do jednego miejsca dziesiętnego).
Kombinatoryczny dowód nierówności $w(n) < 2^n - 1$
dla każdego $n \geq 3$ pozostaje otwarty
(Sekcja~\ref{sec:open}, punkt~(i)).
\end{remark}

\subsection*{4.2. Algorytm inkrementalny}

\begin{lemma}[Aktualizacja okna]
\label{lem:incremental}
Zbiór okien aktywnych dla $n+1$ generatorów wyznaczają
nowe punkty podziału wynikające z $\pi_{n+1}(s) = \pi_n(s) - \pi_n(s - p_{n+1})$
(Lemat~\ref{lem:parity-recurrence}).
Przejście od $n$ do $n+1$ generatorów kosztuje $O(\sigma_n)$.
\end{lemma}

\begin{proof}
Nośnik $\pi_n$ zawiera się w $[0, \sigma_n]$, więc
aktualizacja $\pi_{n+1}(s) = \pi_n(s) - \pi_n(s - p_{n+1})$
wymaga co najwyżej $\sigma_n + 1$ operacji.
\end{proof}

\begin{proof}[Dowód Twierdzenia~\ref{thm:general}]
Wyznaczamy $\mathcal{W}$ --- zbiór par $(lo, hi, c)$ reprezentujących
okna aktywne z wartościami $K^{(n)}$ --- algorytmem inkrementalnym:
zaczynamy od $n = 1$ (jedno okno $[0, p_1)$ z wartością $+1$)
i wykonujemy $n-1$ kroków z Lematu~\ref{lem:incremental}.
Całkowity koszt: $\sum_{k=1}^{n-1} O(\sigma_k) = O(n \cdot \sigma_n)$.

Następnie dla każdego okna $(lo, hi, c) \in \mathcal{W}$
wykonujemy jeden przebieg przesuwnego okna po $G$ w czasie $O(N)$:
akumulujemy $c \cdot \sum_{j \in G \cap [lo, hi)} n_j$.
Łączny koszt: $O(T_{\mathrm{prep}} + w(n) \cdot N)$.
\end{proof}

\begin{remark}
\label{rem:activation}
Zmiana $K^{(n)}$ przy dodaniu generatora $p_{n+1}$ ma postać
\[
  K^{(n+1)}(d) - K^{(n)}(d) =
  \sum_{\substack{S \subseteq \{p_1,\ldots,p_n\} \\ \sigma(S)+p_{n+1} \leq d}}
  (-1)^{|S|+1}.
\]
Pozycja nieaktywna pod $K^{(n)}$ staje się aktywna pod $K^{(n+1)}$
wtedy i tylko wtedy, gdy ta suma jest niezerowa.
Empirycznie $K^{(n)}(d) \notin \{-1, 0, 1\}$ jest możliwe;
wartości $\pm 3$ obserwowane przy $n = 12$
(Tabela~\ref{tab:positions}).
\end{remark}

% ==================================================
\section{Dowód Twierdzenia~\ref{thm:density}}
\label{sec:proof-density}
% ==================================================

\begin{proof}[Dowód Twierdzenia~\ref{thm:density}]
Dowodzimy, że $w(n)/((2^n-1) P(n)) \to 0$.

\textbf{Krok 1.}
Z Lematu~\ref{lem:w-bound}: $w(n) \leq \sigma_n \leq n\,p_n$.

\textbf{Krok 2.}
Ponieważ $p_1 < p_2 < \cdots < p_n$ są parami różne i $\geq 2$,
to $p_k \geq k+1$ dla każdego $k$, skąd
\[
  H_n := \sum_{k=1}^n \frac{1}{p_k} \leq \sum_{j=2}^{n+1} \frac{1}{j} \leq \log(n+1).
\]
Z nierówności $\log(1-x) \geq -2x$ dla $x \in [0, 1/2]$:
\[
  \log P(n) = \sum_{k=1}^n \log\!\Bigl(1 - \tfrac{1}{p_k}\Bigr)
  \geq -2H_n \geq -2\log(n+1),
\]
skąd $P(n) \geq (n+1)^{-2}$.

\textbf{Krok 3.}
Z Kroków 1--2 i $2^n - 1 \geq 2^{n-1}$:
\[
  \frac{w(n)}{(2^n-1)\,P(n)} \leq \frac{2\,n\,p_n\,(n+1)^2}{2^n}.
\]

\textbf{Krok 4.}
Z założenia $\log p_n = o(n)$:
\[
  \log\frac{p_n}{2^n} = \log p_n - n\log 2 = o(n) - n\log 2 \to -\infty,
\]
więc $p_n / 2^n \to 0$, a czynnik $2n(n+1)^2$ rośnie wielomianowo.
Zatem wyrażenie z Kroku~3 dąży do $0$:
czynnik $p_n/2^n = e^{\log p_n - n\log 2} \to 0$
(bo $\log p_n - n\log 2 = o(n) - n\log 2 \to -\infty$),
a wielomian $2n(n+1)^2$ rośnie wolniej niż $2^n/p_n \to \infty$.
Jest to równoważne $\delta(n)/P(n) \to 1$.
\end{proof}

\begin{remark}
\label{rem:density-empirical}
Dla $n \geq 25$ w Tabeli~\ref{tab:mertens}:
$w(n)/(2^n-1) < 10^{-4}$, tzn.\ $w(n)/(2^n-1)$ jest
nierozróżnialne od zera w precyzji czterech miejsc dziesiętnych
$P(n)$ we wszystkich ośmiu testowanych rodzinach ciągów generatorów.
\end{remark}

% ==================================================
\section{Problemy otwarte}
\label{sec:open}
% ==================================================

\begin{enumerate}[(i)]

\item \textbf{Ograniczoność $|K^{(n)}(d)|$.}
Czy $\sup_{n,d} |K^{(n)}(d)| < \infty$?
Empirycznie $|K^{(n)}(d)| \leq 3$ dla $n \leq 12$
(Tabela~\ref{tab:positions}).

\item \textbf{Stabilność wartości.}
Czy dla każdego ustalonego $d$ istnieje $n_0 = n_0(d)$ takie,
że $K^{(n)}(d) = K^{(n_0)}(d)$ dla wszystkich $n \geq n_0$?

\item \textbf{Ostra nierówność $w(n) < 2^n - 1$.}
Twierdzenie~\ref{thm:density} dowodzi $w(n)/(2^n-1) \to 0$,
lecz nie daje ostrego ograniczenia.
Czy $w(n) < 2^n - 1$ dla każdego $n \geq 3$
i dowolnego ciągu wzajemnie pierwszych generatorów?

\item \textbf{Charakteryzacja $n^*$.}
Niech $n^* = \min\{n : \delta(n) > 0\}$.
Twierdzenie~\ref{thm:density} dowodzi $\delta(n)/P(n) \to 1$,
lecz nie charakteryzuje $n^*$.
Hipoteza~\ref{conj:harmonic} sugeruje, że $n^*$ zależy wyłącznie
od sumy harmonicznej $H_n$.

\item \textbf{Monotoniczne zachowanie $\delta(n)$.}
Empirycznie $\delta(n)$ osiąga maksimum przy $n = 16$
i następnie maleje (Tabela~\ref{tab:mertens}).
Czy $\delta(n)$ osiąga właściwe maksimum przy skończonym $n$?

\item \textbf{Struktura $\Delta w(n)$.}
Od $n = 19$: przyrosty $\Delta w(n) = w(n) - w(n-1)$
układają się w dwa rosnące ciągi według parzystości $n$.
Jaki jest mechanizm arytmetyczny tego podziału?

\item \textbf{Asymptotyka $\Delta w(n)$.}
Czy $\Delta w(n) / p_n \to 1$ dla ciągów z $p_n \to \infty$?
Dane numeryczne: $p_{300} = 1987$, $\Delta w(300) = 1986$
(ciąg bazowy); $p_{300} = 2153$, $\Delta w(300) = 2153$
(ciąg od $101$). Zbieżność $\Delta w(n) \sim p_n$ pociągałaby
$w(n) \sim \sigma_n$.

\end{enumerate}

\begin{conjecture}
\label{conj:harmonic}
Niech $n^* = \min\{n : \delta(n) > 0\}$.
Wartość $n^*$ zależy wyłącznie od sumy harmonicznej
$H_n = \sum_{k=1}^n 1/p_k$: ciągi o tej samej $H_n$
co pierwsze $n$ liczb pierwszych mają ten sam $n^*$,
niezależnie od arytmetycznej natury generatorów.
\end{conjecture}

Hipoteza jest potwierdzona eksperymentalnie
(Tabela~\ref{tab:sequences}).
Formalny dowód wymaga charakteryzacji $n^*$
opisanej w problemie~(iv).

\medskip
\textbf{Powiązana literatura.}
Kasowanie parzystościowe we włączeniach--wyłączeniach
pojawia się w niezwiązanych kontekstach:
w teorii wielomianów chromatycznych
(Whitney~\cite{Whitney1932}; Chen i Qian~\cite{ChenQian2018})
oraz w teorii sit analitycznych
(Tao~\cite{Tao2007}).
Rzadkość przedziałowa $K^{(n)}$ i kasowanie
parzystościowe dla jąder półgrup numerycznych
nie były, według naszej wiedzy, wcześniej opisane.

% ==================================================
\section{Wnioski}
\label{sec:conclusion}
% ==================================================

Punktem wyjścia była obserwacja Huanga~\cite{Huang2026}:
forma kwadratowa $Q$ na zbiorze luk półgrupy numerycznej
$\langle a, b \rangle$ odtwarza statystykę kombinatoryczną $\mathrm{dinv}$.
Niniejsza praca odpowiada na pytanie o efektywność obliczeniową $Q$
i ujawnia przy tym strukturę, która nie była wcześniej opisana.

\textbf{Redukcja złożoności.}
Przedziałowa budowa jądra $K$ (Lemat~\ref{lem:K-intervals})
pozwala zastąpić podwójną sumę $O(N^2)$ przez dwa przebiegi
przesuwnego okna w czasie $O(N)$ (Twierdzenie~\ref{thm:linear}).
Przy $n$ generatorach jądro $K^{(n)}$ jest zdefiniowane przez
$2^n$ składników zasady włączeń--wyłączeń, lecz wskutek kasowania parzystościowego posiada co najwyżej
$w(n) \leq \sigma_n$ okien aktywnych (Twierdzenie~\ref{thm:general}).
Empirycznie $w(n) \sim \sigma_n$, podczas gdy $w(n) \ll 2^n$:
wykładnicza kombinatoryka redukuje się do rzadkiej struktury
przedziałowej, a algorytm inkrementalny wyznacza tę strukturę
w czasie $O(n \cdot \sigma_n)$.

\textbf{Asymptotyka.}
Twierdzenie~\ref{thm:density} formalizuje to zjawisko ilościowo:
przy warunku subeksponencjalnym $\log p_n = o(n)$ zachodzi
\[
  \frac{w(n)}{2^n - 1} = o\!\left(\prod_{i=1}^n
    \Bigl(1 - \frac{1}{p_i}\Bigr)\right), \qquad n \to \infty.
\]
Gęstość okien aktywnych jest asymptotycznie zaniedbywalnie mała
względem iloczynu gęstości $P(n)$; równoważnie $\delta(n)/P(n) \to 1$.
Dane numeryczne (Tabela~\ref{tab:mertens}) potwierdzają zbieżność
przy $n \leq 30$ i wskazują na maksimum $\delta(n)$ przy $n = 16$,
którego kombinatoryczna charakteryzacja pozostaje otwarta.

\textbf{Perspektywy.}
Kasowanie parzystościowe w $K^{(n)}$ można ująć algebraicznie jako
działanie operatora różnicowego
$\Delta_{p} f(d) = f(d) - f(d - p)$
na funkcji wskaźnikowej $\mathbf{1}_{[0,\infty)}$:
\[
  K^{(n)} = \Delta_{p_1} \circ \cdots \circ \Delta_{p_n}\,
             \mathbf{1}_{[0,\infty)}.
\]
W tym ujęciu $\pi_n(s)$ jest współczynnikiem maski operatorowej,
a rzadkość $K^{(n)}$ odpowiada zerowym współczynnikom tej maski.
Związek z funkcjami tworzącymi, algebrą operatorową
i analizą harmoniczną na $\mathbb{Z}$ wyznacza kierunek
dalszych badań.

% ==================================================
\appendix
\section{Weryfikacja numeryczna}
\label{sec:verification}
% ==================================================


\section*{A.1. Warunki eksperymentu}
\label{subsec:a1-warunki}

Eksperymenty przeprowadzono w Python~3.14 (SymPy~1.14, NumPy)
na stacji AMD Ryzen~9~5900X (64\,GB, Windows~11~Pro)
i laptopie Intel Core i7-13620H (16\,GB, Windows~11~Pro).
Arytmetyka: IEEE~754 podwójna precyzja ($\varepsilon \approx 2{,}2 \times 10^{-16}$).

Trzy implementacje --- naiwna $O(N^2)$, sumy prefiksowe $O(N\log N)$
i przesuwne okno $O(N)$ --- tworzą łańcuch weryfikacji.
Liczby operacji są deterministycznymi funkcjami $(a, b)$.
Sprawdzenie krzyżowe przez bezpośrednią enumerację $O(2^n)$
przy $n = 27, 28$ przeprowadzono wielowątkowo;
wszystkie pozostałe wartości całkowite
($w(n)$, liczby operacji, pozycje aktywne)
były identyczne bitowo na obu platformach.
Wszystkie skrypty weryfikacyjne dostępne są w repozytorium
\url{https://github.com/Flaman55/RelMathApps/tree/main/number_theory/numerical_semigroups/linear_Q}:
\texttt{verify\_linear\_Q.py} (Twierdzenia~\ref{thm:linear}--\ref{thm:general},
Tabele~\ref{tab:two-gen}--\ref{tab:multi-gen}),
\texttt{appendix\_verify.py} (Tabele~\ref{tab:mertens}--\ref{tab:sequences}),
\texttt{incremental\_windows.py} (algorytm inkrementalny),
oraz \texttt{section\_C\_parallel.py} (obliczenia równoległe, $n=19,\ldots,30$).


\section*{A.2. Dwa generatory}
\label{subsec:a2-dwa}

\begin{table}[H]
\centering
\begin{tabular}{rrrrrr}
\toprule
$a$ & $b$ & $N$ & oper.\ $O(N^2)$ & oper.\ $O(N)$ & maks.\ błąd \\
\midrule
  3 &   5 &     4 &          16 &        16 & $5{,}3\times10^{-15}$ \\
  5 &   7 &    12 &         144 &        48 & $1{,}8\times10^{-14}$ \\
  7 &  11 &    31 &         961 &       124 & $1{,}1\times10^{-13}$ \\
 11 &  13 &    60 &        3600 &       240 & $2{,}8\times10^{-13}$ \\
 23 &  29 &   314 &       98596 &      1256 & $1{,}2\times10^{-11}$ \\
 37 &  41 &   721 &      519841 &      2884 & $4{,}3\times10^{-11}$ \\
 97 & 101 &  4801 &   23049601 &     19204 & $1{,}3\times10^{-9}$  \\
\bottomrule
\multicolumn{6}{l}{\footnotesize Wszystkie 450 par $(a,b)$ z $b \leq 40$: \textbf{PASS}.
  Tolerancja $10^{-9}$ ($10^{-8}$ dla $N > 500$).}
\end{tabular}
\caption{Weryfikacja Twierdzenia~\ref{thm:linear}: dwa generatory.
Stosunek operacji $N^2/(4N) = N/4$.}
\label{tab:two-gen}
\end{table}


\section*{A.3. Wiele generatorów}
\label{subsec:a3-wiele}

\begin{table}[H]
\centering
\begin{tabular}{lrrrr}
\toprule
Generatory & $N$ & $w(n)$ & maks.\ błąd & Status \\
\midrule
$\{2,3\}$     & 1 & 2 & $0$                   & OK \\
$\{3,5\}$     & 4 & 2 & $8{,}9\times10^{-16}$ & OK \\
$\{2,3,5\}$   & 1 & 4 & $0$                   & OK \\
$\{3,5,7\}$   & 3 & 5 & $1{,}8\times10^{-15}$ & OK \\
$\{2,3,5,7\}$ & 1 & 8 & $0$                   & OK \\
\bottomrule
\end{tabular}
\caption{Weryfikacja Twierdzenia~\ref{thm:general}: wiele generatorów
wobec naiwnego $O(N^2)$.}
\label{tab:multi-gen}
\end{table}

\newpage
\section*{A.4. Dane gęstości (Twierdzenie~\ref{thm:density})}
\label{subsec:a4-gestosc}

\begin{table}[H]
\centering
\small
\begin{tabular}{r|rrrrr}
\toprule
$n$ & $2^n-1$ & $w(n)$ & $P(n)$ & Rzadkość\% & $\delta(n)$ \\
\midrule
 2 &          3 &    2 & 0{,}3333 &  0{,}0\% & $-0{,}3333$ \\
 3 &          7 &    4 & 0{,}2667 & 14{,}3\% & $-0{,}3048$ \\
 4 &         15 &    8 & 0{,}2286 & 26{,}7\% & $-0{,}3048$ \\
 5 &         31 &   14 & 0{,}2078 & 29{,}0\% & $-0{,}2438$ \\
 6 &         63 &   20 & 0{,}1918 & 44{,}4\% & $-0{,}1257$ \\
 7 &        127 &   30 & 0{,}1805 & 59{,}1\% & $-0{,}0557$ \\
 8 &        255 &   38 & 0{,}1710 & 72{,}2\% & $+0{,}0220$ \\
 9 &        511 &   48 & 0{,}1636 & 81{,}6\% & $+0{,}0697$ \\
10 &       1023 &   60 & 0{,}1579 & 88{,}0\% & $+0{,}0993$ \\
11 &       2047 &   84 & 0{,}1529 & 92{,}5\% & $+0{,}1118$ \\
12 &       4095 &  116 & 0{,}1487 & 95{,}3\% & $+0{,}1204$ \\
13 &       8191 &  146 & 0{,}1451 & 97{,}2\% & $+0{,}1273$ \\
14 &      16383 &  198 & 0{,}1417 & 98{,}3\% & $+0{,}1296$ \\
15 &      32767 &  236 & 0{,}1387 & 99{,}0\% & $+0{,}1315$ \\
16 &      65535 &  288 & 0{,}1361 & 99{,}4\% & $\mathbf{+0{,}1317}$ \\
17 &     131071 &  352 & 0{,}1338 & 99{,}7\% & $+0{,}1311$ \\
18 &     262143 &  412 & 0{,}1316 & 99{,}8\% & $+0{,}1300$ \\
19 &     524287 &  492 & 0{,}1296 & 99{,}9\% & $+0{,}1287$ \\
20 &    1048575 &  564 & 0{,}1278 & 99{,}9\% & $+0{,}1273$ \\
21 &    2097151 &  646 & 0{,}1260 &100{,}0\% & $+0{,}1257$ \\
22 &    4194303 &  720 & 0{,}1245 &100{,}0\% & $+0{,}1243$ \\
23 &    8388607 &  808 & 0{,}1230 &100{,}0\% & $+0{,}1229$ \\
24 &   16777215 &  892 & 0{,}1216 &100{,}0\% & $+0{,}1215$ \\
25 &   33554431 &  996 & 0{,}1203 &100{,}0\% & $+0{,}1203$ \\
26 &   67108863 & 1096 & 0{,}1191 &100{,}0\% & $+0{,}1191$ \\
27 &  134217727 & 1200 & 0{,}1180 &100{,}0\% & $+0{,}1180$ \\
28 &  268435455 & 1306 & 0{,}1169 &100{,}0\% & $+0{,}1169$ \\
29 &  536870911 & 1416 & 0{,}1158 &100{,}0\% & $+0{,}1158$ \\
30 & 1073741823 & 1528 & 0{,}1148 &100{,}0\% & $+0{,}1148$ \\
\bottomrule
\end{tabular}
\caption{Liczba okien aktywnych $w(n)$, iloczyn gęstości $P(n)$,
rzadkość $1 - w(n)/\sigma_n$
i deficyt $\delta(n)$ (Twierdzenie~\ref{thm:density})
dla $p_i = i$-ta liczba pierwsza.
Wartości graniczne: $n=8$ --- pierwsze $\delta(n) > 0$;
$n=21$ --- rzadkość osiąga $100\%$ (po zaokrągleniu);
$n=25$ --- $w(n)/(2^n-1) < 10^{-4}$.
Maksimum $\delta(n)$ przy $n=16$ (pogrubione).
Wartości ,,100,0\%'' wynikają z zaokrąglenia do jednego miejsca dziesiętnego.
Dane $n \leq 18$: enumeracja bezpośrednia;
$n = 19, \ldots, 30$: algorytm inkrementalny,
niezależnie potwierdzony dla $n=27,28$ enumeracją $O(2^n)$.}
\label{tab:mertens}
\end{table}

\newpage
\section*{A.5. Wartości $K^{(n)}(d)$ --- aktywacja i dezaktywacja}
\label{subsec:a5-wartosci}

\begin{table}[H]
\centering
\small
\begin{tabular}{r|rrrrrrrrrrr}
\toprule
$d$ & $n{=}2$ & $n{=}3$ & $n{=}4$ & $n{=}5$ & $n{=}6$ &
      $n{=}7$ & $n{=}8$ & $n{=}9$ & $n{=}10$ & $n{=}11$ & $n{=}12$ \\
\midrule
 0 &  1 &  1 &  1 &  1 &  1 &  1 &  1 &  1 &  1 &  1 &  1 \\
 3 & -1 & -1 & -1 & -1 & -1 & -1 & -1 & -1 & -1 & -1 & -1 \\
 5 &  $\cdot$ & -1 & -1 & -1 & -1 & -1 & -1 & -1 & -1 & -1 & -1 \\
 7 &  $\cdot$ &  $\cdot$ & -1 & -1 & -1 & -1 & -1 & -1 & -1 & -1 & -1 \\
 8 &  $\cdot$ &  1 &  $\cdot$ &  $\cdot$ &  $\cdot$ &  $\cdot$ &  $\cdot$ &  $\cdot$ &  $\cdot$ &  $\cdot$ &  $\cdot$ \\
 9 &  $\cdot$ &  $\cdot$ &  1 &  1 &  1 &  1 &  1 &  1 &  1 &  1 &  1 \\
10 &  $\cdot$ &  $\cdot$ &  1 &  1 &  1 &  1 &  1 &  1 &  1 &  1 &  1 \\
16 &  $\cdot$ &  $\cdot$ &  $\cdot$ &  $\cdot$ &  1 &  1 &  1 &  1 &  1 &  1 &  1 \\
17 &  $\cdot$ &  $\cdot$ &  $\cdot$ &  1 &  2 &  1 &  1 &  1 &  1 &  1 &  1 \\
25 &  $\cdot$ &  $\cdot$ &  $\cdot$ &  $\cdot$ &  $\cdot$ &  $\cdot$ &  1 &  1 &  1 &  1 &  1 \\
47 &  $\cdot$ &  $\cdot$ &  $\cdot$ &  $\cdot$ &  $\cdot$ &  1 &  1 &  $\cdot$ & -1 & -2 & -3 \\
\bottomrule
\end{tabular}
\caption{Wybrane wartości $K^{(n)}(d)$ dla $n = 2, \ldots, 12$
($p_i = i$-ta liczba pierwsza). Kropka ($\cdot$): $K^{(n)}(d) = 0$.
Pozycja $d = 8$: aktywna przy $n=3$, dezaktywowana przy $n=4$.
Pozycja $d = 17$: $K^{(6)}(17) = 2$.
Pozycja $d = 47$: $K^{(12)}(47) = -3$.
Tabela dostarcza empirycznych przykładów wartości
$|K^{(n)}(d)| = 2, 3$, motywujących problem otwarty~(i).
Zjawisko aktywacji i dezaktywacji pozycji,
empirycznie wynikające z kasowania parzystościowego
między nakładającymi się reprezentacjami sum podzbiorów,
jest treścią problemu otwartego~(ii).}
\label{tab:positions}
\end{table}


\section*{A.6. Weryfikacja Hipotezy~\ref{conj:harmonic}}
\label{subsec:a6-hipoteza}

\begin{table}[H]
\centering
\footnotesize
\begin{tabular}{lrrrr}
\toprule
Ciąg generatorów & Pierw.\ gen. & $n^*$
  & $\delta/P$, $n=25$ & $\delta/P$, $n=300$ \\
\midrule
\multicolumn{5}{l}{\textit{Ciągi pierwsze}} \\
Pierwsze $n$ liczb pierwszych &    2 & 8 & 0{,}9998 & 1{,}0000 \\
Liczby pierwsze od $101$      &  101 & 2 & 0{,}9999 & 1{,}0000 \\
Liczby pierwsze od $1009$     & 1009 & 2 & 0{,}9997 & 1{,}0000 \\
Co druga liczba pierwsza ($p_2, p_4, \ldots$) &    3 & 8 & 0{,}9998 & 1{,}0000 \\
Co trzecia liczba pierwsza ($p_3, p_6, \ldots$) &  5 & 2 & 0{,}9998 & 1{,}0000 \\
Mniejsze z par bliźniaczych   &    3 & 9 & 0{,}9996 & 1{,}0000 \\
Liczby pierwsze $\equiv 3 \pmod{4}$ &  3 & 8 & 0{,}9998 & 1{,}0000 \\
Liczby pierwsze $\equiv 1 \pmod{4}$ &  5 & 2 & 0{,}9998 & 1{,}0000 \\
\multicolumn{5}{l}{\textit{Ciągi losowe (losowe liczby nieparzyste)}} \\
Ta sama $H_n$ co ciąg bazowy (seed 42)  & {---} & 8 & 0{,}9999 & {---} \\
Ta sama $H_n$ co ciąg bazowy (seed 137) & {---} & 8 & 0{,}9999 & {---} \\
Bez dopasowania $H_n$ (seed 42)         & {---} & 2 & 0{,}9999 & {---} \\
\bottomrule
\end{tabular}
\caption{Dane są zgodne z Twierdzeniem~\ref{thm:density}:
$\delta(n)/P(n) \to 1$ dla wszystkich testowanych ciągów.
Przy $n=300$: $\delta(n)/P(n) = 1{,}0000$ (sześć miejsc dziesiętnych)
dla wszystkich ciągów pierwszych.
Ciągi losowe o tej samej $H_n$ co ciąg bazowy dają identyczne $n^*$,
co jest zgodne z Hipotezą~\ref{conj:harmonic}.}
\label{tab:sequences}
\end{table}
\clearpage
\printbibliography

\end{document}
